\section{Confidential customer inference}\label{confidential-customer-inference}

\subsection{Two execution modes}\label{customer-execution-modes}

An ordinary hosted job encrypts text to an operator-controlled miner key.
That miner sees plaintext and could leak or exploit an unpublished
acquisition announcement, earnings release or funding announcement.
Recipient encryption and contractual promises cannot prevent that
authorized endpoint from copying the input.

For sensitive B work, the proposed product uses attested inference:
the customer encrypts to an ephemeral key generated inside a verified
trusted execution environment (TEE). The operator supplies compute and
holds its model weights, but has no protocol key for the customer's
plaintext. This replaces mandatory public B releases as the proposed
privacy route. Ordinary hosting remains an explicitly identified option
for customers accepting operator access; an attested job must never
silently downgrade to it. Attestation is an execution mode for A or B,
not an additional task or reward pool.

A remains downloadable for local use. Customers may run a voluntarily
published B model locally, but the subnet promises no public copy of
the strongest private B model. Local inference uses the customer's
compute and creates no subnet job or fee. Attested hosting instead
depends on qualified hardware, accepted software measurements and the
customer's verifier. Its guarantees are conditional on those components.

\subsection{Attestation and customer key release}\label{attested-customer-inference}

The design follows the separation between hardware evidence, verification
policy and a relying party's decision to release a secret in the RATS
architecture~\cite{rats}. The customer is the relying party for its
input key. Neither a miner's TEE-enabled flag nor a validator's
assertion alone authorizes decryption. The following is a proposed
application protocol, not an implemented attestation service.

\begin{enumerate}
\tightlist
\item
  The customer selects a signed offer binding miner identity, private
  model version, accepted runner/policy hashes, hardware class, output
  key, limits, price and expiry. The client pins its accepted policy and
  vendor trust roots before processing an untrusted offer; the offer
  cannot choose its own reference measurements.
\item
  The selected instance generates a fresh job-scoped encryption key and
  receipt-signing key inside protected memory. It returns evidence bound
  to the customer's unpredictable nonce, those keys, the offer and
  assignment context, the runtime policy and verified model manifest.
  The customer checks signatures, freshness, endorsements, revocation,
  security-version floors and disabled debug access before sending text.
\item
  The CPU evidence must cover the boot chain and enforcement of the
  complete workload policy. For GPU inference, verify the attached GPU
  evidence and its protected channel to that same confidential VM.
  Pairing evidence from an unrelated healthy GPU or ordinary GPU memory
  does not qualify. Pin the full CPU/GPU/driver/firmware combination.
\item
  After successful verification, the customer encrypts directly to the
  attested job key using the bound HPKE context. No gateway, model-owner
  key broker, validator, fleet-wide recovery key or alternate instance
  receives a customer decryption capability. The instance accepts only
  the bound job and verifies the customer's signature over the output key.
\item
  The measured application decrypts, runs the private model, encrypts the
  result only to the customer, and signs the receipt binding the input,
  output and execution version. It then discards job keys and transient
  state according to the enforced lifecycle. Public settlement sees
  permitted metadata and salted commitments, not text or raw text hashes.
\end{enumerate}

An initial hardware candidate is a confidential VM using AMD SEV-SNP or
Intel TDX with a supported NVIDIA confidential-computing GPU and protected
CPU/GPU transfers. NVIDIA documents specific supported configurations;
qualification must pin one tested combination rather than infer coverage
from a GPU product name~\cite{nvidiacc}. GPU attestation checks device
security state and firmware. Application and model binding require the
measured loader and execution policy described in
Section~\ref{attested-b-models}; a GPU quote by itself supplies neither.
A CPU-only enclave with unprotected GPU offload is insufficient.

\subsection{Application isolation and remaining trust}\label{attested-isolation}

Protect input parsing, tokenization, author-profile construction,
generation, search, postprocessing and result encryption inside the
accepted boundary. The miner must have no guest administration, debugger,
shell, plaintext log, crash dump or command channel into that workload.
Disable outgoing tools, external model APIs, Pangram calls and telemetry.
Load model material before accepting customer secrets, make model state
immutable for the job, and prohibit online training or exporting modified
weights. The approved application exposes only authenticated job input
and customer-encrypted output interfaces. Attesting malicious application
code would preserve its ability to leak.

Use isolated job state, with no prompt/KV cache reuse between customers
and no content-dependent public error details. Clear transient buffers
before reuse; forbid host-readable swap or snapshots and unsafe migration.
Protected caches can still expose information through timing, so cache
isolation is an application requirement beyond memory
encryption~\cite{privatecache}. Audit private model inputs as untrusted
data and restrict their operator graph and execution budget. Malicious
weights can influence outputs or timing; encryption of returned text
does not prove absence of every covert channel.

The initial sensitive-job profile must therefore use a qualified fixed
operator graph, fixed tensor-size and token-budget buckets, padding after
early completion, and one padded response at the offer's fixed release
slot. Keep success and failure envelopes the same public size and expose
no token stream or model-selected completion time. This limits deliberate
leakage through output length and early stopping. It costs compute and
latency; microarchitectural and host scheduling side channels still need
independent evaluation. Adaptive graphs or finer-grained metering require
a separately qualified policy.

The client verifier and its update path are part of the trusted computing
base. Publish its source, reproducible builds, accepted measurements,
vendor roots and security-version policy for independent inspection.
Certify reference-policy changes under the strict validator quorum and
activate them prospectively; the customer must accept the new policy
before releasing another key. An operator-controlled web page must not
receive plaintext before this verification. Pin the verifier locally or
in a customer-controlled application.

Require fresh challenge evidence and current revocation collateral before
key release, with fixed maximum evidence age and job duration in the
accepted policy. A failed, stale or unavailable verification stops the
job before disclosure. Restarts or model/runtime changes invalidate the
instance keys and require fresh attestation and customer authorization.
Disable resumable customer state in the initial profile; refuse new key
releases to revoked versions. A patch or revocation cannot undo a secret
already exposed through a vulnerability.

Hardware vendors, firmware, the accepted runtime and verifier, and the
customer's own device remain trust dependencies. Timing, size buckets,
chosen service, prices and availability remain observable; padding and
bounded scheduling reduce some leakage without proving its absence.
The design does not claim protection against every side channel,
physical attack or future hardware exploit. A miner can refuse service
or withhold ciphertext. Attestation does not prove writing quality,
guarantee availability or solve private fair exchange. Sensitive paid
launch depends on an independent security review and the measured
acceptance criteria in Section~\ref{development-and-launch-criteria}.

\subsection{Envelope, assignment and recovery}\label{assigned-miner-encryption}
\label{assignment-recovery-and-the-privacy-boundary}

Use HPKE with X25519/HKDF-SHA256, HKDF-SHA256 and ChaCha20Poly1305,
plus a separate customer signature; HPKE Base does not authenticate the
sender~\cite{hpke}. Bind protocol, chain, subnet, contract, job and
assignment generation, offer hash, recipient key, execution mode and
policy, model version, output key and expiry. For ordinary hosting,
a hotkey-signed miner key record supplies the recipient binding; that
mode grants the operator plaintext access. In attested mode the key
binding additionally requires the evidence above. No alternative
recipient envelope is supplied to validators or the model owner.

The encrypted payload includes source text, reference works, fitted
profiles, brief, private metadata and commitment openings. Use canonical
encoding and a fresh secret 32-byte salt for each commitment:
\begin{equation}\label{eq:input-commitment}
 C=\operatorname{SHA256}\bigl(
 \mathrm{domain}\mathbin\Vert\mathrm{context}\mathbin\Vert
 \mathrm{salt}\mathbin\Vert\mathrm{length}\mathbin\Vert
 \mathrm{payload}\bigr).
\end{equation}
Raw hashes, filenames, descriptive titles and report URLs stay inside
the encrypted payload. The customer alone receives the result envelope,
including probabilities, explanations and edit notes. In attested mode,
the measured runner supplies it; the host cannot replace its output key.

A retry on another instance or miner requires a new customer-authorized
assignment, successful verification for the selected mode, and a newly
encrypted request. There is no automatic plaintext forwarding, key
escrow or downgrade. Loss of a job key can make that job unrecoverable;
the payment policy handles expiry separately. Customers may provide
evidence to a chosen validator under the next subsection, but no audit
exception supplies automatic customer access.

Public A artifacts may be evaluated inside the protected instance.
A confidential customer job cannot send text to Pangram without a
separate customer-authorized disclosure outside this protocol. The
attested receipt proves no Pangram pass. Authorized subnet benchmarks
have their own recipients and external reports; their success is not a
per-document guarantee for an unseen customer revision.

\subsection{Customer-directed evidence
disclosure}\label{customer-directed-evidence-disclosure}

A customer can ask a specific validator to inspect a job by creating a
new evidence packet encrypted to that validator's authenticated key. The
customer signs a disclosure context binding the job, selected validator
and key version, purpose, scope, and validity window. The packet can
contain source passages, the delivered result, the miner's signed
response, and relevant commitment openings. The customer selects what to
send. Neither the miner nor a gateway can activate this access on the
customer's behalf.

This grants the selected validator access to that packet, with no
decryption capability for other jobs or undisclosed material. The
validator must not forward it to other reviewers or providers without
further customer direction. Expiry ends authorization for new
processing; it cannot erase plaintext the recipient already obtained.
Keep the packet, findings, and openings private; publish only any agreed
decision or commitment needed by the protocol.

The validator checks which supplied evidence opens the original
commitments. A redacted excerpt does not open a commitment to a complete
document. Full openings or a separately specified selective-disclosure
commitment scheme are needed for stronger binding. Incomplete evidence
can support a limited opinion, but cannot establish fidelity of the
whole job.

For detection, that validator may reproduce a published A artifact on
the supplied material within the customer's disclosure scope. This grants
no access to other customer text and does not independently reproduce private B weights. An attested
job's receipt may also be checked within the disclosed scope.

Inspection is advisory by default. A signed decision affects escrow only
if the job's original terms gave that validator or selection policy
authority and both parties accepted those terms. Choosing a reviewer
does not itself change payment rights or extend deadlines. This route
permits voluntary inspection while preserving the default assigned-miner
input boundary.
